\section{Nichtlineare Eigenwertprobleme}
In diesem Abschnitt betrachten wir die Minimierung des Energiefunktionals
\begin{equation}\label{equ:1.1}
I(w) := \frac{1}{2} \cdot \int_U |Dw|^2 \td{x}
\end{equation}
mit $w \in H_0^1(U)$ und $w \equiv 0 $ auf $\partial U$.
Wir minimieren $I$ unter der Nebenbedingung
\begin{equation}\label{equ:1.2}
J(w) := \int_U G(w) \td{x} = 0,
\end{equation}
wobei $G : \R  \to \R$ glatt ist.

%\begin{genericthm}{Korollar}\label{skript:1.1}
%Sei $g := G^\prime$ und $|g(z)| \leq  \tilde{C} \cdot (|z| + 1 )$ für alle $z \in \R$.
%Dann gilt
%\begin{align*}
%|G(z)| \leq C \cdot ( |z|^2 + 1))
%\end{align*}
%für alle $ z \in \R$.
%\end{genericthm}
Der Beweis erfolgt hier durch Anwendung des Fundamentalsatzes der Analysis und kurzen Abschätzungen.
Nun nehmen wir an, dass die offene Menge $U$ beschränkt, zusammenhängend und einen glatten Rand besitzt.
Außerdem setzen wir für $g := G^\prime$ die Abschätzung

\begin{align*}
|g(z)| \leq  C \cdot (|z| + 1 )
\end{align*}
voraus.
Weiter definieren wir 
\begin{align*}
\mathcal{A} := \left\lbrace w \in H_0^1(U) \ | \ J(w)=0 \right\rbrace
\end{align*}
die Menge der zugelassenen Funktionen und gelangen zu unserem ersten Satz.

\begin{genericthm}{Existenz eines eingeschränkten Minimierers}\label{skript:1.1}
Sei $\mathcal{A}$ nichtleer.
Dann existiert ein $u \in \mathcal{A}$, sodass
\begin{align*}
I(u) = \min \limits_{w \in \mathcal{A}} I(w)
\end{align*}
erfüllt ist.
\end{genericthm}

\begin{proof}
Zunächst wählen wir uns eine Folge $u_n$ in $\mathcal{A}$ mit
\begin{align*}
I(u_n) \rightarrow m = \min \limits_{w \in \mathcal{A}} I(w).
\end{align*}
Nun ist $H_0^1(U)$ schwach folgenkompakt, wodurch wir mit
\begin{align*}
u_{n_k} \rightharpoonup u
\end{align*}
eine schwache konvergente Teilfolge in $H_0^1(U)$ mit $I(u) \leq m$ finden.
Durch Anwendung des Kompaktheitssatzes von Rellich-Kontrachov erhalten wir 
\begin{align*}
u_{n_k} \rightarrow u
\end{align*}
in $L^2(U)$.
Der Kompaktheitssatz liefert uns eine kompakte Einbettung von $H_0^1(U)$ in $L^2(U)$
und kompakte Operatoren bilden schwach konvergente Folgen auf konvergente Folgen ab.
Durch punktweise Anwendung des Mittelwertsatzes der Differentialrechnung und $u_n \in \mathcal{A}$ gilt
\begin{align*}
|J(u)| 
= |J(u) - J(u_{n_k})|
\leq \int_U | G(u) - G(u_{n_k})| \td{x}
\leq \int_U | u - u_{n_k}| \cdot | 1 + |u| + |u_{n_k}|| \td{x}
\rightarrow 0
\end{align*}
für $k \rightarrow \infty$.
\end{proof}

\begin{genericthm}{Lagrange Multiplikator}\label{skript:1.3}
Sei $u \in \mathcal{A}$ mit
\begin{equation}\label{equ:1.3}
I ( u) = \min \limits_{w \in \mathcal{A}} I(w).
\end{equation}
Dann existiert ein $\lambda \in \R$, so dass
\begin{equation}\label{equ:1.4}
\int_U Du \cdot Dv \td{x}
= \lambda \cdot \int_U g(u) \cdot v \td{x} 
\end{equation}
für alle $v \in H_0^1(U)$ gilt.
\end{genericthm}

\begin{proof}
Sei $v \in H_0^1(U)$ beliebig.
\begin{enumerate}

\item[\textbf{(1)}]
Angenommen $g(u) \neq 0$ fast überall auf $U$.
Dann existiert ein $w \in H_0^1(U)$, so dass
\begin{equation}\label{equ:1.5}
\int_U g(u) \cdot w \td{x} \neq 0
\end{equation}
gilt.
Wir definieren nun
\begin{equation}\label{equ:1.6}
j(\tau,\sigma) 
:= J(u + \tau \cdot v \sigma \cdot w)
\end{equation}
und erhalten direkt $j(0,0) = 0$.
Da $j$ in $C^1$ ist gilt außerdem
\begin{equation}\label{equ:1.7}
\frac{\partial j}{\partial \tau}(\tau,\sigma)
= 
\int_U g( u + \tau v + \sigma w) \cdot v \td{x}
\end{equation}
und
\begin{equation}\label{equ:1.8}
\frac{\partial j}{\partial \sigma}(\tau,\sigma)
= 
\int_U g( u + \tau v + \sigma w) \cdot w \td{x}.
\end{equation} 
Mit (\ref{equ:1.5}) gilt somit 
\begin{equation} \label{equ:1.9}
\frac{\partial j}{\partial \sigma}(\tau,\sigma)(0,0) \neq 0,
\end{equation}
womit nach dem Satz über implizite Funktionen ein $ \phi \in \mathrm{C}^1 $ mit
\begin{equation}\label{equ:1.10}
\phi(0) = 0
\end{equation}
und 
\begin{equation}\label{equ:1.11}
j(\tau, \phi(\tau)) = 0 
\end{equation}
für hinreichend kleine $\tau$ exisitiert.
Durch Differenzieren erhalten wir 
\begin{align*}
\frac{\partial j}{\partial \tau }(\tau, \phi(\tau)) + 
\frac{\partial j}{\partial \sigma}(\tau, \phi(\tau)) \cdot \phi^\prime (\tau) = 0,
\end{align*}
womit mit (\ref{equ:1.7}) und (\ref{equ:1.8})
\begin{equation}\label{equ:1.12}
\phi^\prime(0) 
=
-\frac{\int_U g(u) \cdot v \td{x}}{\int_U g(u) \cdot w \td{x}}
\end{equation}
folgt.
\item[\textbf{(2)}]
Wir setzen nun
\begin{align*}
w(\tau) := \tau v + \phi(\tau) \cdot w
\end{align*}
für hinreichend kleines $\tau$ und definieren
\begin{align*}
i(\tau) := I ( u + w(\tau)).
\end{align*}
Mit (\ref{equ:1.11}) gilt $J(u + w(\tau)) = 0 $, womit $u + w(\tau) \in \mathcal{A}$ gilt.
Die $C^1$ Funktion $i$ besitzt also ein Minimum in der $0$.
Damit erhalten wir
\begin{equation}\label{equ:1.13}
0 = i^\prime(0) = \int_U Du \cdot (Dv + \phi^\prime(0)Dw) \td{x}.
\end{equation}
Mit (\ref{equ:1.12}), (\ref{equ:1.13}) und der Definition
\begin{align*}
\lambda := \frac{\int_U Du \cdot Dw \td{x}}{\int_U g(u)w \td{x}}
\end{align*}
folgt die gewünschte Gleichung
\begin{align*}
\int_U Du \cdot Dv \td{x} = \lambda \cdot \int_U g(u) v \td{x}
\end{align*}
für alle $v \in H_0^1(U)$.

\item[\textbf{(3)}]
Um den Beweis abzuschließen nehmen wir nun an, dass $g(u) = 0$ fast überall gilt.
Wenn wir nun $g$ durch beschränkte Funktionen approximieren, so erhalten wir die schwache Ableitung $DG(u) = g(u) = 0$.
Nun ist $U$ zusammenhängend, womit $G(u)$ konstant fast überall ist.
Mit $J(u) = 0 $ folgt direkt $G(u) = 0 $ fast überall.
Wegen \textit{Theorem 2} in Evans 5.5 gilt  $ u = 0 $ auf $\partial U $ und es folgt $G(0)= 0$.

Also erfüllt $0$ die Nebenbedingung.
Aber dann folgt auch $u = 0$ fast überall, ansonsten würde $I(u) > I(0) = 0$.
Somit folgt unsere Aussage für beliebige $\lambda$.
\end{enumerate}
\end{proof}

\begin{genericdf}{Bemerkung}\label{skript:1.4}
Sei $u$ eine schwache Lösung des nicht-linearen Randwertproblems
\begin{equation}\label{equ:1.14}
\begin{split}
-\Delta u & = \lambda \cdot g(u) \ \text{in} \ U \\
u & = 0 \quad \quad \text{auf} \ \partial U.
\end{split}
\end{equation}
Hierbei ist $\lambda$ der Lagrange-Multiplikator, welcher aus der Nebenbedingung $J(u) = 0$ stammt.
Diesen Zusammenhang können wir bei der Konstruktion des $\lambda$ im letzten Beweis erkennen.
Wir wollen nun die Verbindungen unserer zwei Theoreme und (\ref{equ:1.14}) zeigen.
Da $u$ eine schwache Lösung ist, ist 
\begin{equation*}
\int_U - \Delta u \cdot v  \td{x}
= 
\int_U  Du \cdot Dv \td{x}
=
\lambda \cdot \int_U g(u) v \td{x} 
\end{equation*}
für alle $v \in H_0^1(U)$ erfüllt.
Also ist eine Lösung unseres eingeschränkten Minimierungsproblems automatisch eine Lösung von (\ref{equ:1.14}).
Ein Problem dieser Form mit den Unbekannten $(u,\lambda)$, wobei $u \neq 0 $, nennen
wir \textbf{\textit{nichtlineares Eigenwertproblem}}.  
\end{genericdf}